% Generated draft: requires mathematical and source review.
\documentclass[11pt]{article}
\usepackage[margin=25mm]{geometry}
\usepackage{fontspec}
\setmainfont{Latin Modern Roman}
\usepackage{amsmath,amssymb,mathtools,unicode-math}
\setmathfont{Latin Modern Math}
\usepackage{enumitem,needspace,xurl,hyperref}
\setlength{\parindent}{0pt}
\setlength{\parskip}{5pt}
\setlength{\emergencystretch}{3em}
\newcommand{\N}{\mathbb{N}}
\newcommand{\Z}{\mathbb{Z}}
\newcommand{\Q}{\mathbb{Q}}
\newcommand{\R}{\mathbb{R}}
\newcommand{\C}{\mathbb{C}}
\newcommand{\F}{\mathbb{F}}
\newcommand{\A}{\mathbb{A}}
\newcommand{\PP}{\mathbb{P}}
\newcommand{\E}{\mathbb{E}}
\newcommand{\eps}{\varepsilon}
\newcommand{\dd}{\,\mathrm{d}}
\begin{document}
\section*{UnsolvedMath ID: 30005819}
\textbf{Source problem code:} catalogue entry
\subsection{Definitions and mathematical statement}
Let
\[
\theta_0=\frac{\pi}{2}\bigl(1-10^{-10000}\bigr).
\]
For an integer dimension \(n\geq 1\), let \(X\subseteq\R^n\) be finite. For pairwise distinct points \(x,y,z\in X\), define the angle at \(y\) by
\[
\angle xyz=\arccos\!\left(\frac{\langle x-y,z-y\rangle}{\|x-y\|\,\|z-y\|}\right)\in[0,\pi],
\]
where the inner product and norm are Euclidean. Assume that
\[
\angle xyz\leq\theta_0
\qquad\text{for every ordered triple of pairwise distinct points }x,y,z\in X.
\]
Thus the restriction applies at every vertex of every triangle determined by three distinct points, and equality with \(\theta_0\) is permitted. The proposed assertion is
\[
|X|\leq 1.99^n.
\]
Is this assertion true for every such finite set and every intended dimension?

\emph{Unresolved qualification.} The source does not explicitly specify the range of dimensions or any sufficiently-large-dimension qualification. The convention \(n\geq1\) above makes the ambient dimension explicit, but the intended universal range, including low-dimensional cases and sets for which the triple condition is vacuous, requires review. No asymptotic restriction or exception is being supplied as a correction. This draft follows the supplied reconstructed statement; the raw extraction is incomplete and interleaves neighboring problems.
\subsection{Short English statement}
Must a finite set in Euclidean \(n\)-space have at most \(1.99^n\) points if every angle determined by three distinct points is at most \(\frac{\pi}{2}(1-10^{-10000})\)? The intended dimension range needs clarification. The supplied record labels the question open; this draft does not independently verify that status.
\subsection{Sources}
\begin{itemize}
\item Local registry: \texttt{data/unsolvedmath/problems.json}, supplied record with permanent UnsolvedMath ID \(30005819\), title ``Nearly Right-Angled Point Sets,'' exact source code \texttt{catalogue entry}. The record gives the proposed year as 2024 and supplies no proposer attribution.
\item Supplied primary-source citation: \emph{Discrete Geometry} (2024), Oberwolfach Reports, DOI 10.4171/owr/2024/3, pp. 137--202. \url{https://doi.org/10.4171/owr/2024/3}. The record marks its statement \texttt{corrected\_verified} and says it was reconstructed and checked against the official report. That verification is supplied provenance, not an independent check performed for this draft.
\item Supplied literature source: Kupavskii and Zakharov, \emph{The right acute angles problem?}. \url{https://arxiv.org/abs/1910.00798}.
\item Supplied literature source: \emph{Oberwolfach Report 3/2024}. \url{https://ems.press/content/serial-article-files/48651}.
\item Literature assessment dated exactly \texttt{2026-08-22}: supplied dated curation, not independently verified. It labels the question open and reports that the right acute angles paper gives bounds near \(2^n\) and a lower construction near \((\sqrt{2})^n\), without establishing the proposed \(1.99^n\) bound. The background explicitly calls this dated literature triage derived from the supplied OWR curation and requests independent verification of time-sensitive status and cited claims.
\item The raw extraction also supplies a neighboring reference: K\'arolyi and Solymosi, \emph{Almost Disjoint Triangles in 3-Space}, Discrete \& Computational Geometry 28 (2002), pp. 577--583. It appears in the fragment labeled Problem 15 (Eran Nevo), not as a reference or attribution for the angle question. The extraction subsequently begins an incomplete Problem 17 attributed to J\'anos Pach.
\end{itemize}
\end{document}
